On considère $2n$ personnes distinctes. On souhaite les répartir en $n$
binômes, sans distinguer :
\begin{itemize}
  \item l'ordre des deux personnes à l'intérieur d'un binôme ;
  \item l'ordre dans lequel les binômes sont présentés.
\end{itemize}
Par exemple, avec six personnes,
\[
\{\text{Soa}, \text{Tovo}, \text{Fara}, \text{Hery}, \text{Mamy}, \text{Voara}\},
\]
une partition possible est
\[
\{\{\text{Soa}, \text{Tovo}\}, \{\text{Fara}, \text{Hery}\}, \{\text{Mamy}, \text{Voara}\}\}.
\]
On numérote maintenant les personnes de $1$ à $2n$.
\begin{enumerate}
  \item Pour $n = 2$, écrire toutes les façons de partager les quatre
  personnes $\{1, 2, 3, 4\}$ en deux binômes.
  \item Pour $n = 3$, construire quelques partitions des six personnes
  $\{1, 2, 3, 4, 5, 6\}$ en trois binômes.
  \item À une permutation $(a_1, a_2, \ldots, a_{2n})$, on associe les binômes
  consécutifs
  \[
  \{a_1, a_2\}, \{a_3, a_4\}, \ldots, \{a_{2n-1}, a_{2n}\}.
  \]
  Expliquer pourquoi cette correspondance n'est pas encore une bijection.
  \item Pour une partition fixée, expliquer pourquoi on peut :
  \begin{itemize}
    \item échanger les deux personnes dans chacun des $n$ binômes de $2^n$
    façons ;
    \item permuter les $n$ binômes de $n!$ façons.
  \end{itemize}
  \item En déduire qu'une même partition est représentée par exactement
  \[
  2^n n!
  \]
  permutations.
  \item Comme il existe $(2n)!$ permutations des $2n$ personnes, montrer que
  le nombre de partitions en $n$ binômes est
  \[
  \frac{(2n)!}{2^n n!}.
  \]
  \item Vérifier que, pour six personnes, le nombre de partitions en trois
  binômes est
  \[
  \frac{6!}{2^3\, 3!} = 15.
  \]
  \item On appelle $P_n$ le nombre de partitions de $2n$ personnes en $n$
  binômes. Montrer également que
  \[
  P_n = (2n - 1) P_{n-1}.
  \]
  \item Expliquer cette dernière formule en choisissant d'abord le partenaire
  de la personne numéro $1$.
\end{enumerate}
